\section{Localized Fast Rates, Influence Functions, and Gauge-Singular Fisher Information}
\label{app:local-statistics}
\subsection{Whitened Newton decrement and a leading-constant sharp risk bound}
\label{app:sharp-statistics}
Write $F(\theta)=R(\psi(\theta))$, $F_n(\theta)=R_S(\psi(\theta))$
in a fixed $D$-dimensional regular chart with $\nabla F(0)=0$.
For a positive definite reference matrix $H$ define
$\|\theta\|_H^2=\theta^\top H\theta$ and
$d_n^2=\nabla F_n(0)^\top H^{-1}\nabla F_n(0)$.
The reference $H$ is independent of the sample in the concentration
statement, or is fixed after conditioning on a noise-independent design.

\begin{theorem}[Whitened local risk certificate]
\label{thm:sharp-newton-risk}
Suppose on $\|\theta\|_H\leq r$ that
\begin{align}
(1-\delta_H)H&\preceq\nabla^2F_n(\theta),&
\nabla^2F(\theta)&\preceq(1+\epsilon_H)H,
\qquad 0\leq\delta_H<1 .
\end{align}
If $d_n<(1-\delta_H)r$, the constrained empirical minimum is unique
and interior, and
\begin{align}
\|\widehat\theta_n\|_H&\leq\frac{d_n}{1-\delta_H},\\
F(\widehat\theta_n)-F(0)
&\leq\frac{1+\epsilon_H}{2(1-\delta_H)^2}d_n^2 .
\label{eq:sharp-newton-risk}
\end{align}
For an empirical $\epsilon_T$-minimizer in the same ball,
\begin{equation}
F(\theta_T)-F(0)
\leq\frac{1+\epsilon_H}{2}
\left\{\frac{d_n}{1-\delta_H}
+\sqrt{\frac{2\epsilon_T}{1-\delta_H}}\right\}^{\!2}.
\label{eq:sharp-optimization-statistics}
\end{equation}
\end{theorem}
\begin{proof}
Use $z=H^{1/2}\theta$. The empirical Hessian is at least
$(1-\delta_H)I$ and its center gradient has norm $d_n$.
On $\|z\|=r$ its outward derivative is at least
$(1-\delta_H)r^2-d_nr>0$. Thus the minimum is interior.
Strong monotonicity gives
$(1-\delta_H)\|\widehat z\|\leq d_n$.
The population Hessian bound yields
$F(\theta)-F(0)\leq(1+\epsilon_H)\|\theta\|_H^2/2$.
Empirical strong convexity bounds
$\|\theta_T-\widehat\theta_n\|_H^2
\leq2\epsilon_T/(1-\delta_H)$.
Apply the triangle inequality and keep the square unexpanded.
For the quadratic pair
\begin{align}
F(\theta)-F(0)&=\tfrac12(1+\epsilon_H)\|\theta\|_H^2,\\
F_n(\theta)-F_n(0)&=s^\top\theta+
\tfrac12(1-\delta_H)\|\theta\|_H^2,
\end{align}
the minimizer is $-(1-\delta_H)^{-1}H^{-1}s$ and
Eq.~\eqref{eq:sharp-newton-risk} is an equality.
Thus both the coefficient $1/2$ and the relative-Hessian factor
$(1-\delta_H)^{-2}$ are sharp for these stated sector conditions.
\end{proof}

\begin{lemma}[Laurent--Massart and Hsu--Kakade--Zhang score concentration]
\label{lem:sharp-quadratic-tail}
Suppose the independent centered whitened scores
$w_i=H^{-1/2}\nabla\ell_0(z_i)$ satisfy
\begin{equation}
\E e^{u^\top w_i}\leq e^{u^\top\mathsf V u/2},
\qquad u\in\R^D,\quad\mathsf V\succeq0 .
\end{equation}
Then, for every $u_0>0$, with probability at least $1-e^{-u_0}$,
\begin{equation}
d_n^2\leq\frac1n
\left\{\operatorname{tr}\mathsf V
+2\sqrt{\operatorname{tr}(\mathsf V^2)u_0}
+2\|\mathsf V\|_2u_0\right\}.
\label{eq:sharp-score-tail}
\end{equation}
\end{lemma}
\begin{proof}
This is the Gaussian quadratic-form bound of
Laurent--Massart \cite{laurent2000adaptive} and its sub-Gaussian
extension \cite{hsu2012tail}; the short proof also fixes the constants.
Put $W=n^{-1/2}\sum_iw_i$ and let $g\sim N(0,I_D)$ be independent.
For $0<t<(2\|\mathsf V\|_2)^{-1}$,
\begin{align}
\E e^{t\|W\|^2}
&=\E_g\E_W e^{\sqrt{2t}g^\top W}
\leq\E_g e^{t g^\top\mathsf Vg}
=\det(I-2t\mathsf V)^{-1/2},\\
\log\E e^{t(\|W\|^2-\operatorname{tr}\mathsf V)}
&\leq\frac{t^2\operatorname{tr}(\mathsf V^2)}
{1-2t\|\mathsf V\|_2}.
\end{align}
The scalar inequality used here is
$-\log(1-x)-x\leq x^2/[2(1-x)]$, $0\leq x<1$.
Chernoff optimization with
\begin{equation}
t=\frac{\sqrt{u_0}}
{\sqrt{\operatorname{tr}(\mathsf V^2)}
+2\|\mathsf V\|_2\sqrt{u_0}}
\end{equation}
gives the tail threshold in Eq.~\eqref{eq:sharp-score-tail};
$W=\sqrt n\,H^{-1/2}\nabla F_n(0)$.
The zero-covariance case follows by continuity.
\end{proof}

\begin{corollary}[A condition-number-free fast rate]
\label{cor:sharp-risk-tail}
Suppose the Hessian event in Theorem~\ref{thm:sharp-newton-risk}
fails with probability at most $q_H$ and its radius condition is
satisfied by the score threshold. With probability at least
$1-q_H-e^{-u_0}$,
\begin{align}
F(\widehat\theta_n)-F(0)
&\leq\frac{1+\epsilon_H}{2n(1-\delta_H)^2}
\left\{\operatorname{tr}\mathsf V
+2\sqrt{\operatorname{tr}(\mathsf V^2)u_0}
+2\|\mathsf V\|_2u_0\right\}.
\label{eq:sharp-risk-tail}
\end{align}
In the isotropic whitened case $\mathsf V=\sigma^2I_D$ this becomes
\begin{equation}
\frac{(1+\epsilon_H)\sigma^2}
{2n(1-\delta_H)^2}\{D+2\sqrt{Du_0}+2u_0\}.
\label{eq:sharp-isotropic-tail}
\end{equation}
\end{corollary}
\begin{proof}
Intersect the score-tail event with the Hessian event and substitute
Eq.~\eqref{eq:sharp-score-tail}. No independence of these two events
is needed. Whitening retains the covariance--curvature interaction;
there is no separate $\Lambda/\mu^2$ prefactor.
\end{proof}

\begin{proposition}[A Newton--Kantorovich localization certificate]
\label{prop:sharp-kantorovich}
Let $A=\nabla^2F_n(0)\succ0$ and $d=\|A^{-1/2}\nabla F_n(0)\|$.
Suppose, on an $A$-ball containing the radius below,
\begin{equation}
\|A^{-1/2}\{\nabla^2F_n(\theta)-A\}A^{-1/2}\|_2
\leq K\|\theta\|_A .
\end{equation}
If $2Kd<1$, define
\begin{equation}
r_-=\frac{1-\sqrt{1-2Kd}}{K}
=\frac{2d}{1+\sqrt{1-2Kd}},
\qquad d\leq r_-\leq2d .
\label{eq:sharp-kantorovich}
\end{equation}
Then the local stationary point is unique in this ball and
$\|\widehat\theta_n\|_A\leq r_-$.
If $\nabla^2F(\theta)\preceq(1+\epsilon_A)A$ there,
\begin{equation}
F(\widehat\theta_n)-F(0)\leq\tfrac12(1+\epsilon_A)r_-^2 .
\end{equation}
For $K=0$, interpret $r_-=d$.
\end{proposition}
\begin{proof}
This is the affine-invariant Newton--Kantorovich majorant argument
\cite{deuflhard2004newton}, proved here through Banach's contraction
principle. In whitened coordinates consider the frozen-Jacobian map
\begin{equation}
\mathcal N(\theta)=\theta-A^{-1}\nabla F_n(\theta).
\end{equation}
Taylor's integral remainder and the derivative premise give
\begin{align}
\|\mathcal N(\theta)\|_A
&\leq d+\tfrac K2\|\theta\|_A^2,\\
\|D\mathcal N(\theta)\|_{A\to A}
&\leq K\|\theta\|_A .
\end{align}
The majorant equation $d+Kr_-^2/2=r_-$ makes the ball invariant,
and $Kr_-=1-\sqrt{1-2Kd}<1$ makes it a contraction.
Its fixed point solves $\nabla F_n=0$ and the Hessian is positive
throughout the ball. The population-risk bound follows by integration.
The argument is a localization certificate; it does not require
changing QuotientOpt into a Newton iteration.
\end{proof}


\subsection{An explicit finite-sample event construction}
\begin{assumption}
\label{ass:local-statistical-chart}
Let $\psi:\Theta\to\mathcal M_r$ be a fixed regular rank chart,
$\Theta=\{\theta\in\R^D:\|\theta\|_2\leq R_0\}$,
$D=r(m+d-r)$, and set
\begin{align}
\ell_\theta(z)&=\ell(\psi(\theta);z),&
F(\theta)&=\E\ell_\theta(z),&
F_n(\theta)&=\frac1n\sum_{i=1}^n\ell_\theta(z_i),\\
\nabla F(0)&=0,&
2\mu I&\preceq\nabla^2F(\theta)\preceq\Lambda I,\quad\theta\in\Theta .
\end{align}
The observations $z_i$ are IID. Per-example coordinate Hessians satisfy
\begin{align}
\|\nabla^2\ell_\theta(z)\|_2&\leq B,&
\|\nabla^2\ell_\theta(z)-\nabla^2\ell_{\theta'}(z)\|_2
&\leq M\|\theta-\theta'\|_2 .
\end{align}
The centered score $s(z)=\nabla\ell_0(z)$ has directional
sub-Gaussian constant $\nu$:
\begin{equation}
\E e^{u^\top s(z)}\leq e^{\nu^2\|u\|^2/2},
\qquad u\in\R^D .
\label{eq:local-score-model}
\end{equation}
The model uses its stated derivative and moment conditions; a
bounded-loss assumption is not needed in this section.
\end{assumption}

\begin{lemma}[A uniform empirical Hessian event with explicit constants]
\label{lem:local-Hessian-concentration}
For $0<q<1$, the sample-size condition
\begin{equation}
n\geq\frac{32B^2}{\mu^2}
\left[D\log9+
D\log\left(1+\frac{8MR_0}{\mu}\right)+\log\frac4q\right]
\label{eq:local-Hessian-sample}
\end{equation}
implies, with probability at least $1-q/2$,
\begin{equation}
\sup_{\theta\in\Theta}
\|\nabla^2F_n(\theta)-\nabla^2F(\theta)\|_2\leq\mu,
\qquad \nabla^2F_n(\theta)\succeq\mu I .
\end{equation}
\end{lemma}
\begin{proof}
For a symmetric matrix $A$ and a $1/4$ net $\mathcal V$ of the unit
sphere,
\begin{align}
\|A\|_2&\leq
\max_{v\in\mathcal V}|v^\top Av|+\tfrac12\|A\|_2,\\
\|A\|_2&\leq2\max_{v\in\mathcal V}|v^\top Av|,&
|\mathcal V|&\leq9^D .
\end{align}
Take a $\delta_\theta=\mu/(4M)$ net of $\Theta$ of cardinality
at most $(1+8MR_0/\mu)^D$, with centers in $\Theta$.
For $M=0$, a single coordinate center suffices.
For fixed $\theta,v$, the IID variables
$v^\top\nabla^2\ell_\theta(z_i)v$ lie in $[-B,B]$, so
\begin{equation}
\Pr\left\{\left|v^\top
(\nabla^2F_n(\theta)-\nabla^2F(\theta))v\right|>\mu/4\right\}
\leq2e^{-n\mu^2/(32B^2)} .
\end{equation}
A union bound gives operator error at most $\mu/2$ at all net
centers. Off the net, the empirical and population Hessians each
have Lipschitz constant $M$, adding at most
$2M\delta_\theta=\mu/2$. The sample-size condition bounds the union
probability by $q/2$.
\end{proof}

\begin{lemma}[Score concentration in the intrinsic dimension]
With probability at least $1-q/2$,
\begin{equation}
\|\bar s_n\|_2
=\left\|\frac1n\sum_i s(z_i)\right\|_2
\leq t_n
:=\nu\sqrt{\frac{8(D\log5+\log(2/q))}{n}} .
\label{eq:local-score-concentration}
\end{equation}
\end{lemma}
\begin{proof}
A $1/2$ sphere net $\mathcal V$ has at most $5^D$ elements and
$\|a\|\leq2\max_{v\in\mathcal V}v^\top a$.
By the sub-Gaussian premise and exponential Markov inequality,
\begin{align}
\Pr\{v^\top\bar s_n>u\}&\leq e^{-nu^2/(2\nu^2)},\\
\Pr\{\|\bar s_n\|>t\}&\leq5^De^{-nt^2/(8\nu^2)} .
\end{align}
Set $t=t_n$.
\end{proof}

\begin{theorem}[A localized $1/n$ excess-risk bound]
\label{thm:local-fast-rate}
Suppose Eq.~\eqref{eq:local-Hessian-sample} holds and
$t_n\leq\mu R_0/2$. With probability at least $1-q$, the minimum
$\widehat\theta_n=\arg\min_{\theta\in\Theta}F_n(\theta)$ is unique,
interior, and satisfies
\begin{align}
\nabla F_n(\widehat\theta_n)&=0,&
\|\widehat\theta_n\|_2&\leq t_n/\mu,\\
F(\widehat\theta_n)-F(0)
&\leq\frac{\Lambda t_n^2}{2\mu^2}
=\frac{4\Lambda\nu^2}{\mu^2 n}
\{D\log5+\log(2/q)\}.
\label{eq:local-fast-rate}
\end{align}
More generally, every $\theta_T\in\Theta$ with
$F_n(\theta_T)-F_n(\widehat\theta_n)\leq\epsilon_T$ obeys
\begin{equation}
F(\theta_T)-F(0)
\leq\frac{\Lambda}{2}
\left(\frac{t_n}{\mu}+\sqrt{\frac{2\epsilon_T}{\mu}}\right)^2 .
\label{eq:local-optimization-statistics}
\end{equation}
\end{theorem}
\begin{proof}
Work on the intersection of the two preceding events.
For $\|\theta\|=R_0$,
\begin{align}
\theta^\top\nabla F_n(\theta)
&=\theta^\top\bar s_n+
\int_0^1\theta^\top\nabla^2F_n(t\theta)\theta\,dt\\
&\geq-R_0t_n+\mu R_0^2>0 .
\end{align}
Thus a boundary point is improved by moving inward.
Compactness gives an interior minimizer, and strong convexity gives
uniqueness. Strong monotonicity implies
\begin{align}
\mu\|\widehat\theta_n\|^2
&\leq\langle\nabla F_n(\widehat\theta_n)-\nabla F_n(0),
\widehat\theta_n\rangle\\
&=-\langle\bar s_n,\widehat\theta_n\rangle
\leq t_n\|\widehat\theta_n\| .
\end{align}
Population smoothness and $\nabla F(0)=0$ give
$F(\theta)-F(0)\leq\Lambda\|\theta\|^2/2$.
Finally,
\begin{align}
\tfrac\mu2\|\theta_T-\widehat\theta_n\|^2
&\leq F_n(\theta_T)-F_n(\widehat\theta_n)\leq\epsilon_T,\\
\|\theta_T\|&\leq
\|\widehat\theta_n\|+\|\theta_T-\widehat\theta_n\| .
\end{align}
Combining the inequalities proves both claims.
\end{proof}

\begin{corollary}[Closing the deterministic optimization--statistics chain]
If the empirical quotient iteration is in the certified basin of
Theorem~\ref{thm:local-nonlinear-basin} around its empirical minimum,
one may insert
\begin{equation}
\epsilon_T\leq C_{\rm init}q_{\mathsf P}^{2T}
\end{equation}
into Eq.~\eqref{eq:local-optimization-statistics}.
In particular, $\epsilon_T\leq t_n^2/(2\mu)$ is sufficient for
\begin{equation}
F(\theta_T)-F(0)\leq 2\Lambda t_n^2/\mu^2 .
\end{equation}
For $C_{\rm init}>t_n^2/(2\mu)$ this holds whenever
\begin{equation}
T\geq
\frac{\log(2\mu C_{\rm init}/t_n^2)}
{2\log(1/q_{\mathsf P})}.
\label{eq:local-statistical-stopping}
\end{equation}
\end{corollary}
\begin{proof}
The optimization theorem gives an empirical suboptimality bound.
Then $\sqrt{2\epsilon_T/\mu}\leq t_n/\mu$ in
Eq.~\eqref{eq:local-optimization-statistics}. Solving the geometric
inequality gives the required iteration count.
The sample size $n$ counts independent observations. Repeated
optimization epochs on those observations affect $\epsilon_T$,
not $n$.
\end{proof}

\begin{proposition}[Obtaining the coordinate curvature premise]
Let a normal chart at $X_\star$ have
$\mathsf H=\nabla^2F(0)\succeq h_\star I$, and suppose
$\nabla^2F$ is $K_F$-Lipschitz there. If
$R_0\leq h_\star/(2K_F)$, then
\begin{equation}
\nabla^2F(\theta)\succeq h_\star I/2=2\mu I,
\qquad \mu=h_\star/4 .
\end{equation}
At an exact RIP sensing solution one may take
$h_\star\geq a(1-\delta)$.
\end{proposition}
\begin{proof}
Weyl's inequality gives
$\lambda_{\min}(\nabla^2F(\theta))
\geq h_\star-K_F\|\theta\|$.
At the chart center, coordinate and quotient Hessians agree in the
orthonormal normal frame. The RIP Hessian result at $E=0$ supplies
the last bound. A bound on the actual coordinate Hessian variation
is still needed to choose $R_0$.
\end{proof}

\subsection{Influence functions and the exact leading generalization constant}
The local estimating-equation expansion is the standard influence-function
argument for regular M-estimators \cite{vandervaart1998asymptotic}; the
quantitative remainder below makes its dependence on the rank chart explicit.
\begin{theorem}[A quantitative influence expansion]
\label{thm:local-influence}
Write
\begin{align}
\mathsf H&=\nabla^2F(0)\succ0,&
\mathsf H_n&=\nabla^2F_n(0),&
\mathsf C&=\E[s(z)s(z)^\top].
\end{align}
On an event where the empirical Hessian is at least $\mu I$,
its Lipschitz constant is at most $M$, and the local minimum is
interior,
\begin{align}
\widehat\theta_n&=-\mathsf H^{-1}\bar s_n+\Delta_n,\\
\|\Delta_n\|
&\leq\|\mathsf H^{-1}\|_2
\left\{\frac{\|\mathsf H_n-\mathsf H\|_2\|\bar s_n\|}{\mu}
+\frac{M\|\bar s_n\|^2}{2\mu^2}\right\}.
\label{eq:local-influence}
\end{align}
If the event has probability tending to one,
$\mathsf H_n-\mathsf H=o_p(1)$, and the score has finite covariance,
\begin{align}
\sqrt n\,\widehat\theta_n
&\xrightarrow{\ d\ }N(0,\mathsf H^{-1}\mathsf C\mathsf H^{-1}),\\
n\{F(\widehat\theta_n)-F(0)\}
&\xrightarrow{\ d\ }
\tfrac12 Z^\top\mathsf H^{-1}Z,\qquad Z\sim N(0,\mathsf C).
\label{eq:local-influence-clt}
\end{align}
\end{theorem}
\begin{proof}
The stationarity equation and a Taylor remainder give
\begin{align}
0&=\bar s_n+\mathsf H_n\widehat\theta_n+\rho_n,&
\|\rho_n\|&\leq\tfrac M2\|\widehat\theta_n\|^2,\\
\widehat\theta_n+\mathsf H^{-1}\bar s_n
&=\mathsf H^{-1}
\{(\mathsf H-\mathsf H_n)\widehat\theta_n-\rho_n\}.
\end{align}
Insert $\|\widehat\theta_n\|\leq\|\bar s_n\|/\mu$.
The IID score central limit theorem gives
$\sqrt n\,\bar s_n\Rightarrow N(0,\mathsf C)$, and the displayed
error is $o_p(n^{-1/2})$.
The population Taylor expansion then proves the quadratic limit.
\end{proof}

\begin{theorem}[Leading expected risk, training optimism, and generalization]
\label{thm:local-optimism}
Under Assumption~\ref{ass:local-statistical-chart}, consider the
localized empirical minimum equal to $0$ when the uniform Hessian
and interior-score event fails. Assume $\E|\ell_0(z)|<\infty$.
For fixed chart and dimension, as $n\to\infty$,
\begin{align}
\E[F(\widehat\theta_n)-F(0)]
&=\frac{\operatorname{tr}(\mathsf C\mathsf H^{-1})}{2n}
+O(n^{-3/2}),\\
\E[F_n(\widehat\theta_n)-F_n(0)]
&=-\frac{\operatorname{tr}(\mathsf C\mathsf H^{-1})}{2n}
+O(n^{-3/2}),\\
\E[F(\widehat\theta_n)-F_n(\widehat\theta_n)]
&=\frac{\operatorname{tr}(\mathsf C\mathsf H^{-1})}{n}
+O(n^{-3/2}).
\label{eq:local-optimism}
\end{align}
\end{theorem}
\begin{proof}
The two concentration arguments imply a failure probability
$O(e^{-cn})$ for a fixed sufficiently small score threshold and
the fixed chart. Independence and sub-Gaussian scores yield
\begin{align}
\E[\bar s_n\bar s_n^\top]&=\mathsf C/n,&
\E\|\bar s_n\|^4&=O(n^{-2}),&
\E\|\mathsf H_n-\mathsf H\|_2^2&=O(n^{-1}).
\end{align}
The last bound follows by bounding operator norm by Frobenius norm
and summing the entrywise variances of the empirical Hessian.
On the successful event, the influence expansion gives
\begin{align}
F(\widehat\theta_n)-F(0)
&=\tfrac12\bar s_n^\top\mathsf H^{-1}\bar s_n+\mathcal E_n,\\
F_n(\widehat\theta_n)-F_n(0)
&=-\tfrac12\bar s_n^\top\mathsf H^{-1}\bar s_n+\widetilde{\mathcal E}_n,\\
|\mathcal E_n|+|\widetilde{\mathcal E}_n|
&\leq C\{\|\mathsf H_n-\mathsf H\|_2\|\bar s_n\|^2+
\|\bar s_n\|^3\}.
\end{align}
Hessian norms are bounded, and the influence bound justifies
absorbing higher-order terms into this inequality on the
successful event. Cauchy--Schwarz gives expectation $O(n^{-3/2})$.
Off that event the localized differences are zero; the omitted
quadratic score expectation is at most
\begin{equation}
C(\E\|\bar s_n\|^4)^{1/2}
\Pr(\text{failure})^{1/2}=O(n^{-1}e^{-cn/2}).
\end{equation}
Taking the trace proves the first two statements. For the third,
use the exact decomposition
\begin{align}
F(\widehat\theta_n)-F_n(\widehat\theta_n)
&=[F(\widehat\theta_n)-F(0)]
-[F_n(\widehat\theta_n)-F_n(0)]\\
&\quad+F(0)-F_n(0),
\end{align}
whose last term has mean zero.
\end{proof}

\subsection{Isotropic sensing and a sharp intrinsic dimension constant}
\begin{definition}
Let
\begin{align}
y&=\langle Z,X_\star\rangle+\varepsilon,&
\E[\varepsilon\mid Z]&=0,&
\E[\varepsilon^2\mid Z]&=\sigma_\varepsilon^2,\\
\ell(X;(Z,y))&=\tfrac12(\langle Z,X\rangle-y)^2,&
\E[\operatorname{vec}Z(\operatorname{vec}Z)^\top]&=I_{md}.
\label{eq:local-isotropic-model}
\end{align}
All required local derivatives and moments are assumed integrable.
Bounded-feature conditions for the finite-sample Hessian theorem
and finite-moment conditions for asymptotic influence results are
separate premises.
\end{definition}

\begin{theorem}[The intrinsic efficiency constant]
\label{thm:local-isotropic-efficiency}
For $J=D\psi(0)$, represented as an $md\times D$ matrix,
\begin{align}
F(\theta)-F(0)
&=\tfrac12\|\psi(\theta)-X_\star\|_F^2,\\
s(z)&=-\varepsilon J^\top\operatorname{vec}Z,&
\mathsf H&=J^\top J,&
\mathsf C&=\sigma_\varepsilon^2\mathsf H,\\
\operatorname{tr}(\mathsf C\mathsf H^{-1})
&=\sigma_\varepsilon^2D
=\sigma_\varepsilon^2r(m+d-r).
\label{eq:local-intrinsic-efficiency}
\end{align}
Whenever the expectation premises of Theorem~\ref{thm:local-optimism}
or Theorem~\ref{thm:noise-nonlinear-average} apply, their leading
excess-risk constant is therefore $\sigma_\varepsilon^2D/2$.
\end{theorem}
\begin{proof}
The conditional mean-zero noise removes the population cross term:
\begin{align}
\E[\ell(X;z)-\ell(X_\star;z)]
&=\tfrac12\E\langle Z,X-X_\star\rangle^2
=\tfrac12\|X-X_\star\|_F^2 .
\end{align}
At the exact solution the population residual gradient vanishes,
so the coordinate Hessian is $J^\top J$. The conditional noise
variance gives
\begin{align}
\E[s(z)s(z)^\top]
&=\E[\varepsilon^2J^\top\operatorname{vec}Z
(\operatorname{vec}Z)^\top J]\\
&=\sigma_\varepsilon^2J^\top J .
\end{align}
The chart is regular, so $J$ has rank $D$ and $\mathsf H$ is
invertible. The trace identity follows.
\end{proof}

\begin{proposition}[Rank-one Gaussian measurements are isotropic]
For independent $a\sim N(0,I_m)$ and $b\sim N(0,I_d)$ with
$Z=ab^\top$,
\begin{equation}
\E[Z_{ij}Z_{k\ell}]
=\E[a_ia_k]\E[b_jb_\ell]
=\delta_{ik}\delta_{j\ell}.
\end{equation}
Thus Eq.~\eqref{eq:local-isotropic-model} applies to these
measurements under the stated noise and local moment premises.
It does not require independent entries in $Z$.
\end{proposition}

\subsection{Fisher degeneracy along gauge fibers}
\begin{theorem}[Full factor Fisher information and its identifiable image]
\label{thm:local-fisher-gauge}
Assume $\varepsilon\mid Z\sim N(0,\sigma_\varepsilon^2)$ and
isotropic features. At a regular factor pair $p_\star=(B_\star,A_\star)$,
let $L_\star=D\pi_{p_\star}$ with the Euclidean factor inner product.
The per-observation Fisher information is
\begin{align}
\mathsf I_{\rm fac}
&=\sigma_\varepsilon^{-2}L_\star^*L_\star,&
\ker\mathsf I_{\rm fac}&=\mathcal V_{p_\star},&
\operatorname{rank}\mathsf I_{\rm fac}&=D,\\
L_\star\mathsf I_{\rm fac}^\dagger L_\star^*
&=\sigma_\varepsilon^2P_{X_\star}.
\label{eq:local-fisher-gauge}
\end{align}
Hence the identifiable product-space covariance has trace
$\sigma_\varepsilon^2D$, independent of the $r^2$ factor gauge
coordinates.
\end{theorem}
\begin{proof}
The conditional log-likelihood score in a factor direction $\xi$ is
\begin{align}
D\log p(y\mid Z,p_\star)[\xi]
&=\sigma_\varepsilon^{-2}\varepsilon\langle Z,L_\star\xi\rangle .
\end{align}
Taking the expected product of two directional scores gives
$\sigma_\varepsilon^{-2}\langle L_\star\xi,L_\star\zeta\rangle$.
The kernel and rank follow from Theorem~\ref{thm:quotient-space}.
For an SVD $L_\star=\mathsf U\mathsf S\mathsf V^\top$ on its
nonzero singular subspaces,
\begin{align}
L_\star(L_\star^*L_\star)^\dagger L_\star^*
&=\mathsf U\mathsf S(\mathsf S^2)^{-1}\mathsf S\mathsf U^\top\\
&=\mathsf U\mathsf U^\top=P_{X_\star}.
\end{align}
The factor $\sigma_\varepsilon^2$ proves the last identity.
\end{proof}


\begin{corollary}[Design-adaptive efficiency by Fisher whitening]
\label{cor:sharp-general-design}
Replace isotropy by a feature covariance
$\mathsf M=\E[\operatorname{vec}Z(\operatorname{vec}Z)^\top]$
that is positive definite on $T_{X_\star}\mathcal M_r$.
For the same correctly specified homoscedastic model,
\begin{align}
\mathsf H&=J^\top\mathsf M J,&
\mathsf C&=\sigma_\varepsilon^2J^\top\mathsf M J
=\sigma_\varepsilon^2\mathsf H,\\
\operatorname{tr}(\mathsf C\mathsf H^{-1})
&=\sigma_\varepsilon^2D,\\
\mathsf I_{\rm fac}
&=\sigma_\varepsilon^{-2}L_\star^*\mathsf M L_\star,\\
L_\star\mathsf I_{\rm fac}^\dagger L_\star^*
&=\sigma_\varepsilon^2
(P_{X_\star}\mathsf M P_{X_\star})^\dagger,\\
\operatorname{tr}\{\mathsf M
L_\star\mathsf I_{\rm fac}^\dagger L_\star^*\}
&=\sigma_\varepsilon^2D .
\label{eq:sharp-design-cancellation}
\end{align}
Consequently the leading prediction-risk constant
$\sigma_\varepsilon^2D/2$ holds for every such identifiable design
under the stated asymptotic expectation premises.
\end{corollary}
\begin{proof}
The conditional variance identity gives
$\mathsf C=\sigma_\varepsilon^2\mathsf H$ directly.
For $L_\star=\mathsf U\mathsf S\mathsf V^\top$ on the nonzero
singular subspaces, let
$\mathsf M_T=\mathsf U^\top\mathsf M\mathsf U\succ0$. Then
\begin{align}
(L_\star^*\mathsf M L_\star)^\dagger
&=\mathsf V\mathsf S^{-1}\mathsf M_T^{-1}
\mathsf S^{-1}\mathsf V^\top,\\
L_\star(L_\star^*\mathsf M L_\star)^\dagger L_\star^*
&=\mathsf U\mathsf M_T^{-1}\mathsf U^\top
=(P_{X_\star}\mathsf M P_{X_\star})^\dagger .
\end{align}
Taking the $\mathsf M$-weighted trace yields $D$.
This cancellation concerns prediction risk. Product Frobenius error
instead has leading trace
$\sigma_\varepsilon^2\operatorname{tr}(\mathsf M_T^{-1})$,
which retains the design conditioning.
\end{proof}

\begin{proposition}[A Cram\'er--Rao bound for locally unbiased chart estimators]
In the preceding Gaussian model, the Cram\'er--Rao information inequality
\cite{vandervaart1998asymptotic} gives the following result. At a fixed chart parameter let
$\widehat\theta$ be unbiased in a neighborhood and satisfy the usual
differentiation-under-integral conditions. Then
\begin{align}
\operatorname{Cov}(\widehat\theta)
&\succeq\frac{\sigma_\varepsilon^2}{n}\mathsf H^{-1},\\
\E\left[\tfrac12(\widehat\theta-\theta)^\top
\mathsf H(\widehat\theta-\theta)\right]
&\geq\frac{\sigma_\varepsilon^2D}{2n}.
\label{eq:local-cramer-rao}
\end{align}
\end{proposition}
\begin{proof}
Let $S_n$ be the likelihood score for $n$ independent observations.
At the fixed parameter,
\begin{align}
\E S_n&=0,&
\E[S_nS_n^\top]&=n\mathsf I,&
\mathsf I&=\sigma_\varepsilon^{-2}\mathsf H .
\end{align}
Differentiating $\E_\theta\widehat\theta=\theta$ gives
$\E[(\widehat\theta-\theta)S_n^\top]=I_D$.
For $A=(n\mathsf I)^{-1}$,
\begin{align}
0&\preceq
\E[(\widehat\theta-\theta-AS_n)
(\widehat\theta-\theta-AS_n)^\top]\\
&=\operatorname{Cov}(\widehat\theta)-(n\mathsf I)^{-1}.
\end{align}
Take the $\mathsf H/2$-weighted trace. This lower bound concerns
the stated locally unbiased class and quadratic tangent loss.
\end{proof}
